%! The Tangent Function — Unit 7, Lesson 7.1 — Group Activity (Answer Key)
\documentclass[10pt]{article}
\usepackage{precalculus-article}
\usepackage{precalculus-key}   % pulls in -boxes; do NOT also load -boxes
\newcommand{\TallMath}[1]{$\displaystyle #1\rule[-1.4em]{0pt}{3.2em}$}

\begin{document}

\pageheader{Unit 7, Lesson 7.1}{Group Activity --- Answer Key}
\namepartnerperiod

\vspace{0.08in}

\begin{scenariobox}[The Tangent Function --- Tiered Practice]{plumlight}
The tangent function connects the unit circle to slope and to periodic asymptotic
behavior.  Choose \textbf{one tier} below based on your readiness level.
Everyone should be prepared to explain their reasoning to the class.
\end{scenariobox}

\vspace{0.06in}

%----------------------------------------------------------------------
% Tier R
%----------------------------------------------------------------------
\begin{tcolorbox}[colback=white, colframe=black!40,
    title=\textbf{Tier R --- Remediate},
    fonttitle=\bfseries, arc=2mm, left=3mm, right=3mm, top=2mm, bottom=2mm]

\textbf{Context:} The unit circle values below are given.  Use them to compute
$\tan\theta = \sin\theta / \cos\theta$ at each angle.

\vspace{6pt}
\renewcommand{\arraystretch}{1.55}
\begin{tabular}{|c|c|c|c|}
\hline
$\theta$ & $\sin\theta$ & $\cos\theta$ & $\tan\theta$ \\
\hline
$0$          & $0$          & $1$          & \ans{$0$} \\
\hline
$\pi/6$      & $1/2$        & $\sqrt{3}/2$ & \ans{$1/\sqrt{3} = \sqrt{3}/3 \approx 0.577$} \\
\hline
$\pi/4$      & $\sqrt{2}/2$ & $\sqrt{2}/2$ & \ans{$1$} \\
\hline
$\pi/3$      & $\sqrt{3}/2$ & $1/2$        & \ans{$\sqrt{3} \approx 1.732$} \\
\hline
$\pi/2$      & $1$          & $0$          & \ans{undefined} \\
\hline
$3\pi/4$     & $\sqrt{2}/2$ & $-\sqrt{2}/2$& \ans{$-1$} \\
\hline
$\pi$        & $0$          & $-1$         & \ans{$0$} \\
\hline
\end{tabular}

\vspace{8pt}
\textbf{1.}\quad At which angles above is $\tan\theta$ \emph{undefined}?  Why?

\ans{$\tan\theta$ is undefined at $\theta = \pi/2$ because $\cos(\pi/2) = 0$ and
division by zero is undefined.}

\vspace{4pt}
\textbf{2.}\quad The tangent function has vertical asymptotes at $\theta = \pi/2 + k\pi$
for integer $k$.  List the three asymptotes closest to $\theta = 0$.

\ans{$\theta = -\pi/2$, $\theta = \pi/2$, $\theta = 3\pi/2$}

\vspace{4pt}
\textbf{3.}\quad Between $\theta = 0$ and $\theta = \pi/2$, is $\tan\theta$
increasing or decreasing?  How do you know from the table?

\ans{Increasing: the values go $0, \sqrt{3}/3, 1, \sqrt{3}, \ldots$ --- each output is larger than the previous one as $\theta$ increases.}

\end{tcolorbox}

\vspace{0.06in}

%----------------------------------------------------------------------
% Tier A
%----------------------------------------------------------------------
\begin{tcolorbox}[colback=white, colframe=black!40,
    title=\textbf{Tier A --- Approaching Proficiency},
    fonttitle=\bfseries, arc=2mm, left=3mm, right=3mm, top=2mm, bottom=2mm]

\textbf{Context:} The graph below shows one and a half periods of a tangent function
$f(\theta) = a\tan(\theta) + d$.

\vspace{6pt}
\begin{center}
\begin{tikzpicture}[xscale=1.6, yscale=0.45]
  \draw[->] (-1.8,0) -- (4.2,0) node[right]{\small $\theta$};
  \draw[->] (0,-4.5) -- (0,6.5) node[above]{\small $y$};
  \draw[dashed,slate] (-1.571,-4.3)--(-1.571,6.3);
  \draw[dashed,slate] ( 1.571,-4.3)--( 1.571,6.3);
  \draw[dashed,slate] ( 4.712,-4.3)--( 4.712,6.3);
  \draw[navy,thick,domain=-1.45:-0.01,samples=80,smooth]
        plot (\x,{2*tan(deg(\x))+1});
  \draw[navy,thick,domain=0.01:1.45,samples=80,smooth]
        plot (\x,{2*tan(deg(\x))+1});
  \draw[navy,thick,domain=1.70:3.10,samples=80,smooth]
        plot (\x,{2*tan(deg(\x))+1});
  \fill[navy] (0,1) circle (2pt) node[left,font=\scriptsize]{$(0,1)$};
  \fill[navy] (0.785,3) circle (2pt) node[right,font=\scriptsize]{$(\pi/4,3)$};
  \foreach \y in {-4,-2,2,4,6}
    \draw (0.04,\y)--(-0.04,\y) node[left,font=\scriptsize]{$\y$};
  \foreach \x/\l in {-1.571/$-\frac{\pi}{2}$,0.785/$\frac{\pi}{4}$,1.571/$\frac{\pi}{2}$,3.142/$\pi$}
    \draw (\x,0.1)--(\x,-0.1) node[below,font=\scriptsize]{\l};
\end{tikzpicture}
\end{center}

\vspace{4pt}
\textbf{1.}\quad What is the period of $f$?

\ans{$\pi$ (the distance between consecutive asymptotes is $\pi$)}

\vspace{4pt}
\textbf{2.}\quad List all vertical asymptotes visible in the graph.

\ans{$\theta = -\pi/2$, $\theta = \pi/2$, $\theta = 3\pi/2$}

\vspace{4pt}
\textbf{3.}\quad On the interval $(-\pi/2, \pi/2)$, is $f$ concave up or concave down
for $\theta < 0$?  For $\theta > 0$?

\ans{Concave down for $\theta < 0$; concave up for $\theta > 0$.  The curve bends downward (like a hill) to the left of the origin and upward (like a bowl) to the right.}

\vspace{4pt}
\textbf{4.}\quad Find $a$ and $d$.

\ans{At $\theta = 0$: $f(0) = a\tan(0)+d = d = 1$, so $d = 1$.\\[4pt]
At $\theta = \pi/4$: $f(\pi/4) = a\cdot 1 + 1 = 3$, so $a = 2$.\\[4pt]
Therefore $f(\theta) = 2\tan(\theta) + 1$.}

\end{tcolorbox}

\vspace{0.06in}

%----------------------------------------------------------------------
% Tier E
%----------------------------------------------------------------------
\begin{tcolorbox}[colback=white, colframe=black!40,
    title=\textbf{Tier E --- Extension},
    fonttitle=\bfseries, arc=2mm, left=3mm, right=3mm, top=2mm, bottom=2mm]

\textbf{Context:} Analyze $g(\theta) = 2\tan\!\left(\dfrac{\theta}{2} - \dfrac{\pi}{4}\right) + 1$.

\vspace{6pt}
\textbf{1.}\quad Rewrite in the form $a\tan(b(\theta+c))+d$.

\ans{Factor: $\dfrac{\theta}{2} - \dfrac{\pi}{4} = \dfrac{1}{2}\!\left(\theta - \dfrac{\pi}{2}\right)$.\\[4pt]
So $g(\theta) = 2\tan\!\left(\tfrac{1}{2}(\theta - \tfrac{\pi}{2})\right) + 1$.\\[4pt]
$a = 2$, $b = \tfrac{1}{2}$, $c = -\tfrac{\pi}{2}$, $d = 1$.}

\vspace{4pt}
\textbf{2.}\quad Period of $g$:

\ans{Period $= \pi/|b| = \pi\div(1/2) = 2\pi$.}

\vspace{4pt}
\textbf{3.}\quad Vertical asymptotes:

\ans{Set $\tfrac{1}{2}(\theta - \tfrac{\pi}{2}) = \tfrac{\pi}{2} + k\pi$:\\[4pt]
$\theta - \tfrac{\pi}{2} = \pi + 2k\pi$\\[4pt]
$\theta = \tfrac{3\pi}{2} + 2k\pi$ for integer $k$.\\[4pt]
So asymptotes at $\ldots, -\tfrac{\pi}{2}, \tfrac{3\pi}{2}, \tfrac{7\pi}{2}, \ldots$}

\vspace{4pt}
\textbf{4.}\quad Compute $g(0)$ and $g(\pi/2)$:

\ans{$g(0) = 2\tan\!\left(-\tfrac{\pi}{4}\right) + 1 = 2(-1)+1 = -1$.\\[4pt]
$g(\pi/2) = 2\tan\!\left(\tfrac{\pi}{4}-\tfrac{\pi}{4}\right)+1 = 2\tan(0)+1 = 2(0)+1 = 1$.}

\vspace{4pt}
\textbf{5.}\quad Description of transformations:

\ans{Compared to $y=\tan\theta$: the graph is vertically dilated by a factor of 2 (outputs doubled), has period $2\pi$ instead of $\pi$ (horizontal stretch by factor 2), shifted right by $\pi/2$ (phase shift), and shifted up by 1 unit (vertical translation).}

\end{tcolorbox}

\end{document}
